%%%%%%%%%%%%%%%%%%%% preface.tex %%%%%%%%%%%%%%%%%%%%%%%%%%%%%
% Preface for Springer Monograph
%%%%%%%%%%%%%%%%%%%%%%%%%%%%%%%%%%%%%%%%%%%%%%%%%%%%%%%%%%%%%

\preface

For nearly four decades, the computing revolution drew its momentum from the harmonic cadence of Moore's Law~\cite{Moore1965} and Dennard Scaling~\cite{Dennard1974}: transistors shrank, clock frequencies accelerated into the multi-gigahertz regime, and single-threaded execution performance soared without threatening silicon thermal limits. Around 2005, the irreversible breakdown of Dennard scaling and the onset of the Power Wall~\cite{Borkar1999} halted clock frequency scaling, forcing a historic pivot toward spatial parallelism: multi-core processors, massively parallel graphics processors, and heterogeneous Systems-on-Chip (SoCs).

Yet scaling compute elements from tens to hundreds, and now thousands of cores on a single die or across 2.5D/3D multi-chiplet assemblies, immediately collided with an unyielding physical reality: the \textbf{interconnect crisis}~\cite{Ho2001,Dally2001}. Global shared multi-drop buses collapsed under quadratic $RC$ wire delay and serialization contention; full crossbar matrices exploded in $O(N^2)$ silicon area and routing track congestion; and deep-submicron copper interconnects below 14~nm succumbed to severe electron scattering and resistive degradation~\cite{Mayadas1970}.

The \textbf{Network-on-Chip (NoC)} paradigm answered this crisis by transposing packet-switched networking principles directly onto silicon~\cite{Dally2001,Benini2002}. Today, on-chip and die-to-die communication fabrics form the nervous system of modern computing---from mobile application processors and high-performance server SoCs to wafer-scale accelerators and neuromorphic tensor processors.

\section*{Scope and Structure of This Book}
This textbook offers a rigorous, first-principles exploration of on-chip and multi-chiplet interconnection networks. Rather than treating interconnects as abstract graph topologies, this volume bridges the chasm between physical silicon reality, router microarchitecture, mathematical deadlock theory, and emerging technology frontiers.

The book is organized into five cohesive parts:
\begin{itemize}
    \item \textbf{Part I: Foundations \& Router Microarchitecture (Chapters 1--3)} establishes the physical interconnect crisis, mathematical delay formulations, switching techniques, credit and token flow control, and pipelined router microarchitecture with virtual channel allocation and crossbar arbitration.
    \item \textbf{Part II: Topologies, Routing \& Deadlock Freedom (Chapters 4--5)} investigates direct, indirect, and high-radix networks, Dally's channel dependency graph (CDG) theorems, Duato's adaptive routing theory, and turn models.
    \item \textbf{Part III: Physical Constraints: Power, Thermal \& Reliability (Chapters 6--7)} addresses dark silicon dynamics, voltage/frequency scaling, compact thermal RC modeling, hotspot mitigation, and fault-tolerant routing.
    \item \textbf{Part IV: Advanced Packaging, 3D NoC \& Chiplets (Chapters 8--9)} explores through-silicon via (TSV) vertical architectures, 2.5D silicon interposers, and the Universal Chiplet Interconnect Express (UCIe) die-to-die standard.
    \item \textbf{Part V: Emerging Frontiers (Chapter 10)} covers silicon nanophotonic interconnects (ONoC), micro-ring resonators, and on-chip millimeter-wave/sub-THz wireless communication (WNoC).
\end{itemize}

\section*{Grounding in Premier Literature \& Silicon Implementations}
Every theoretical formulation, architectural trade-off, and silicon case study in this monograph is grounded directly in premier peer-reviewed literature from IEEE and ACM venues, including \emph{IEEE Transactions on Very Large Scale Integration (VLSI) Systems} (TVLSI), \emph{IEEE Transactions on Computer-Aided Design of Integrated Circuits and Systems} (TCAD), \emph{IEEE Journal of Solid-State Circuits} (JSSC), and \emph{IEEE Transactions on Circuits and Systems} (TCAS-I/II). Landmark silicon achievements---such as the 2,048-core HammerBlade SoC~\cite{Gao2026}, the FlooNoC open-source AXI4 fabric~\cite{Fischer2025}, the SWIFT token flow-control architecture~\cite{Davidson2023}, and the open UCIe standard~\cite{UCIe2023}---are dissected alongside first-principles physics.

\section*{Co-Design with the Interactive Companion Hub}
To enhance pedagogy and provide hands-on experimentation, this monograph is co-designed with an interactive online platform:
\begin{center}
\textbf{NoC Learning Hub}: \url{https://network-on-chip.github.io/learning/}
\end{center}
Each chapter links directly to interactive browser-based simulators, including the 2D Mesh Routing Visualizer, Arbiter Arena, Credit Commander, and UCIe Protocol Analyzer, accompanied by multi-tier assessment engines and executable Verilog/Python code repositories.

\vspace{1cm}
\begin{flushright}
\textbf{The NoC-Lab Research Consortium}\\
\emph{November 2026}
\end{flushright}
